\documentclass[10pt,a4paper]{amsart}
\usepackage[margin=19mm]{geometry}
\usepackage{amsmath,amssymb,mathtools}
\usepackage[scr=rsfs,cal=euler]{mathalfa}
\usepackage{xfrac}
\usepackage[unicode,hidelinks,bookmarks=false]{hyperref}
\newcommand{\ie}{i.e.\ }
\newcommand{\cf}{cf.\ }
\newcommand{\resp}{respectively\ }
\newcommand{\Subsection}{\subsection}
\providecommand{\nobreakdash}{\nobreak\hbox{-}}
\setlength{\emergencystretch}{2em}
\multlinegap0pt
\usepackage[shortlabels]{enumitem}
\usepackage{csquotes}
\usepackage{float}
\usepackage{amssymb}
\usepackage[normalem]{ulem}
\usepackage{xcolor}
\newtheorem{proposition}{Proposition}
\def \bm {{\bf m}}
 \def \s {{\sigma}}
\title{The Ising Model on a Two-Community Stochastic Block Model}
\author{Alessandra Bianchi, Vanessa Jacquier, Matteo Sfragara}
\date{Source: arXiv:2604.20631v2 (CC BY 4.0)}
\begin{document}
\maketitle

\noindent\emph{Source and licence.} The Ising Model on a Two-Community Stochastic Block Model. Version arXiv:2604.20631v2 (CC BY 4.0), distributed by arXiv under the Creative Commons Attribution 4.0 International licence, http://creativecommons.org/licenses/by/4.0/, as recorded in the arXiv licence metadata for this exact version and shown on the abstract page https://arxiv.org/abs/2604.20631v2. Full licence notice: \texttt{licenses/arxiv-2604.20631-CC-BY-4.0.txt}.\par\smallskip

\noindent\textit{Editorial scope.} Counted unit: the article's proposition Z\_approx (source0.tex lines 956--974). The article's complete original proof of it is delivered with it (source0.tex lines 976--1191, one \texttt{proof} environment of 216 lines). No mathematics is written, repaired or re-derived: the statement and the proof are the article's own lines, in the article's own notation, and no retained line is substituted or re-flowed.  Retained uncounted as the source-local context the proof needs: lines 332--334; lines 559--561; lines 763--768; lines 770--773; lines 913--920.  Nothing was omitted from any retained span, so the package carries no omission record.  No retained line contains a citation, so the wrapper carries no bibliography and no bibliography entry.  No retained line contains a figure environment, an image-inclusion command or drawing code, so no image is delivered and the package declares no asset dependency.  Editorial formatting only, in addition: an AMS \texttt{amsart} preamble that restates the article's own declarations and macros used by the retained lines, namely the environments \texttt{proposition} on separate counters, together with the standard packages those lines need.  The retained environments print in this document as Proposition 1 in order of appearance, whereas the article prints the same items with the numbering of its own full document.  The complete native source of the article as distributed in the arXiv e-print for this exact version is bundled unchanged as \texttt{sources/198956/source0.tex}.
\begin{align}\label{def:hamiltonian_CW_magnet}
\bar{H}_n(\sigma)=-\frac{n}{2} E_{\alpha_n}(\bm^{(n)}(\s)).
\end{align}

\begin{equation}\label{StirBound}
\frac{c}{\sqrt{n}}e^{- \frac{n}{2}\mathcal{I}(m_i)}\le\binom{n/2}{\frac{1+m_i}{2} n/2}\le Ce^{- \frac{n}{2}\mathcal{I}(m_i)},\qquad \mbox{ for } i=1,2.
\end{equation}

\begin{equation}\label{Max}
\begin{split}
&M_1(\alpha)=G_{\alpha,\beta}(\bm^1)\equiv G_{\alpha,\beta}(\bm^2),\\
&M_2(\alpha)=G_{\alpha,\beta}(\bm^3)\equiv G_{\alpha,\beta}(\bm^4).
\end{split}
\end{equation}

\begin{equation}
\label{def:M1M2}
 \varepsilon(\alpha)= M_1(\alpha)-M_2(\alpha).
\end{equation}

\begin{equation} \label{def-intorni}
\begin{split}
\mbox{-\,\, if } \beta< \beta_c, \,\,\,   &B_{0}= \left \{ \bm \in \Gamma_{n/2}^2: \, \|\bm-\bm^0\| \leq \frac{1}{n^\delta} \right \} \,\, \text{and}
    \,\, \mathcal{C}_0=\Gamma_{n/2}^2 \setminus B_0;\\
\mbox{-\,\, if } \beta> \beta_c, \,\,\,    & B_{i}= \left \{ \bm \in \Gamma_{n/2}^2: \, \|\bm-\bm^i\| \leq \frac{1}{n^\delta} \right \},\, i=1,\ldots, k \,\,\text{and}
    \,\, \mathcal{C}=\Gamma_{n/2}^2 \setminus \bigcup_{i=1}^k B_i.  
\end{split}
\end{equation}

\begin{proposition}\label{Z_approx} 
For $n$ large enough, we have
\begin{equation*}
\bar{Z}_{n,\beta}=
\begin{cases}
e^{n\log2} D_0^{(n)}(1+o(1)), & \text{if }\beta<\beta_c, \\[5pt]
e^{\frac{n}{2} M_1(\alpha_n)} \sum_{i=1}^2 D_i^{(n)} (1+o(1)),& \text{if } \beta  \in \big(\beta_c,\beta^*(\alpha_n)\big], \\[5pt]
    e^{\frac{n}{2} M_1(\alpha_n)} \left (\sum_{i=1}^2 D_i^{(n)}+e^{-\frac{n}{2} \varepsilon(\alpha_n)} \sum_{i=3}^4 D_i^{(n)} \right ) (1+o(1)),  & \text{if } \beta>\beta^*(\alpha_n),
\end{cases} 
\end{equation*}
where $\varepsilon(\alpha_n)$ is defined in $\eqref{def:M1M2}$.
Moreover, if $\beta=\beta_c$, we get
\begin{equation}\label{Zcrit-bound}
\sqrt[4]{n} \, e^{n\log2}  D_{\text{low}}^{(n)}\,(1+o(1))\le \bar{Z}_{n,\beta_c}\le
\sqrt{n} \, e^{n\log2}  D_{\text{up}}^{(n)}\,(1+o(1)), 
\end{equation}
where $D_{\text{low}}^{(n)}$ and $D_{\text{up}}^{(n)}$ are explicit 
uniformly bounded function of $n$, converging to finite positive constants.
\end{proposition}

\begin{proof}
We begin by assuming $\beta>\beta_c$, which corresponds to the case of multiple maximizers, with $k\in\{2,4\}$. 
From  \eqref{def:hamiltonian_CW_magnet}, we can write
\begin{align*}
    \bar{Z}_{n,\beta} &=\sum_{\bm \in \Gamma_{n/2}^2} \binom{n/2}{\frac{1+m_1}{2}n/2} \binom{n/2}{\frac{1+m_2}{2}n/2} e^{\frac{\beta n}{2} E_{\alpha_n}(\bm)} = \sum_{i=1}^k \bar{Z}_{n,\beta}(B_i) +  \bar{Z}_{n,\beta}(\mathcal{C}),
\end{align*}
where we used the partition $\Gamma_{n/2}^2=\bigcup_{i=1}^k B_i \cup \mathcal{C}$ and the notation
$$\bar{Z}_{n,\beta}(A)= \sum_{\bm \in A} \sum_{\substack{\s\in\mathcal{X}_n:\\\bm(\s)=\bm}}e^{-\beta \bar{H}_n(\s)}\qquad \forall A\subseteq \Gamma_{n/2}^2\,.$$

We first analyze the contributions from $\bm\in B_i$, and without loss of generality we consider $i=1$. 
In $B_1$, the Stirling formula can be strengthened to the following asymptotical expression
\begin{align}\label{eq-coeffbin}
    \binom{n/2}{\frac{1+m_j}{2} n/2}= \frac{2}{\sqrt{\pi n (1-m_j^2)}}e^{- \frac{n}{2}\mathcal{I}(m_j)}(1+o(1)),
    \quad\mbox{for } j=1,2,
\end{align}
so that
\begin{align}\label{eq-Z-tipica}
    \bar{Z}_{n,\beta}(B_1)=\sum_{\bm \in B_1} \frac{4}{\pi n \sqrt{1-m_1^2} \sqrt{1-m_2^2}}e^{\frac{n}{2}G_{\alpha_n,\beta}(\bm)} (1+o(1)).
\end{align}
By Taylor expanding $G_{\alpha_n,\beta}(\bm)$ around the global maximizer $\bm^1=(m_g,m_g)$, we can write
\begin{align}\label{taylor_G}
    G_{\alpha_n,\beta}(\bm)= & \,\,G_{\alpha_n,\beta}(\bm^1)+\frac{1}{2} \left ( \frac{\beta}{2}-\frac{1}{1-m_g^2} \right )  \| \bm-\bm^1 \|^2+\frac{\beta \alpha_n}{2}(m_1-m_g)(m_2-m_g) \notag \\
    &+ O(\|\bm-\bm^1\|^3).
\end{align}
%
Substituting this term in \eqref{eq-Z-tipica}, and by the change of variable ${\bf{x}}=\sqrt{n}({\bf{m}}-{\bf{m}}^{1})$, we get
\begin{equation}\label{eq-Z-tipica2}
\begin{split}
    \bar{Z}_{n,\beta}(B_1)
    &=e^{\frac{n}{2}G_{\alpha_n,\beta}(\bm^1)}\sum_{{\bf x} \in R_1^{(n)}} \frac{4}{\pi n \sqrt{1- \left (\frac{x_1}{\sqrt{n}}+m_g \right )^2} \sqrt{1- \left (\frac{x_2}{\sqrt{n}}+m_g \right )^2}}  \\
    & \qquad\qquad\qquad\qquad\quad\,\,\, \cdot e^{ \frac{1}{4} \left ( \frac{\beta}{2}-\frac{1}{1-m_g^2} \right ) \|{\bf x}\|^2  +\frac{\beta \alpha_n}{4} x_1x_2}(1+o(1))  \\
    &= e^{\frac{n}{2}M_1(\alpha_n)} D_1^{(n)} (1+o(1)), 
\end{split}
\end{equation}
where we also used that, by definition \eqref{def-intorni}, $n\|\bm-\bm^1\|^3\le n^{1-3\delta}=o(1)$ if $\delta>1/3$. 
In the same way, we can provide analogous estimates of $Z_{n,\beta}(B_i)$, for $i=2,...,k$. 

We conclude by analyzing the contribution to $\bar{Z}_{n,\beta}$ coming from the configurations in $\mathcal{C}$. Suppose $k=4$ and divide $\mathcal{C}$ into the following four disjoint subsets:
\begin{equation}\label{4subset}
\begin{array}{ll}
\mathcal{C}_1= \{ \bm \in \mathcal{C} \, | \, m_1,m_2\ge 0\},\qquad& \mathcal{C}_2= \{ \bm \in \mathcal{C} \, | \, m_1,m_2<0\},\\ 
 \mathcal{C}_3= \{ \bm \in \mathcal{C} \, | \, m_1<0,m_2\ge 0\},\qquad& \mathcal{C}_4= \{ \bm \in \mathcal{C} \, | \, m_1\ge 0,m_2<0\}.
\end{array}
\end{equation}
We already stress that when $k=2$, the last two subsets are excluded from $\mathcal{C}$, but the remainder of the proof proceeds in the same way. We introduce the following notation for the local maxima of $G_{\alpha_n,\beta}$ on the above subsets:
\begin{equation*}
\begin{split}
    & g_1(\bm)=\max_{\bm \in \mathcal{C}_1} G_{\alpha_n,\beta}(\bm), 
    \qquad  g_2(\bm)=\max_{\bm \in \mathcal{C}_2}G_{\alpha_n,\beta}(\bm), \\
    & g_3(\bm)=\max_{\bm \in \mathcal{C}_3}G_{\alpha_n,\beta}(\bm), 
    \qquad g_4(\bm)=\max_{\bm \in \mathcal{C}_4}G_{\alpha_n,\beta}(\bm).
\end{split}
\end{equation*}
By the Stirling bound \eqref{StirBound}, recalling the notation introduced in \eqref{Max}, we get
\begin{align*}
    \bar{Z}_{n,\beta}(\mathcal{C}) 
    & \leq C n^2 \left (e^{\frac{n}{2}M_1(\alpha_n)} \sum_{i=1,2} e^{-\frac{n}{2}(M_1(\alpha_n)-g_i(\bm))}
    +e^{\frac{n}{2}M_2(\alpha_n)}\sum_{i=3,4}e^{-\frac{n}{2}(M_2(\alpha_n)-g_i(\bm))}
        \right ).
\end{align*}
Note that by construction, for all $i=1,\ldots, 4$ and for all $\bm \in \mathcal{C}_i$,
\begin{equation*}
n|M_i(\alpha_n)-g_i(\bm)|>n\|\bm-\bm^i\|^2>n^{1-2\delta}.
\end{equation*}
Hence we obtain 
\begin{align}\label{stat-atipica}
\bar{Z}_{n,\beta}(\mathcal{C})& \leq Cn^2e^{-\frac{1}{2}n^{1-2\delta}} (2e^{\frac{n}{2}M_1(\alpha_n)}+2e^{\frac{n}{2}M_2(\alpha_n)})<e^{\frac{n}{2}M_1(\alpha_n)} o(1)+e^{\frac{n}{2}M_2(\alpha_n)} o(1),
\end{align}
where in the last inequality, we use that $\delta<1/2$.
Together, \eqref{eq-Z-tipica2} and \eqref{stat-atipica} show that
\begin{align*}
        \bar{Z}_{n,\beta}=\left (e^{\frac{n}{2} M_1(\alpha_n)} \sum_{i=1}^2 D_i^{(n)}+\textbf{1}_{\{k>2\}}e^{\frac{n}{2}M_2(\alpha_n)} \sum_{i=3}^4 D_i^{(n)} \right ) (1+o(1)),
    \end{align*}
hence concluding the analysis for $k=2,4$.


A similar argument applies for $k=1$, whenever $\beta<\beta_c$.
In that case, the main contribution to the partition function comes 
from the neighborhood $B_0$ of the unique maximizer $\bm^0$ of $G_{\alpha_n,\beta}$, and the above procedure can be implemented even more directly. %\frac{2}{1+\alpha_n}.

In the critical regime $\beta=\beta_c$, when we still have the unique maximizer $\bm^0$, some extra care is needed to deal with the fact that
the Hessian of $G_{\alpha_n,\beta_c}$ in $\bm^0$  has a null eigenvalue.
In that case, a Taylor expansion around $\bm^0$ up to the fourth order is needed to deal with cases where the small order derivatives vanish as $n\to\infty$, depending on the value of $\alpha_n$. 
Explicitly, we get
\begin{align}\label{eq:Exp-G-Crit}
   G_{\alpha_n,\beta_c}(\mathbf{m})
&= 2\log 2
-\frac{\alpha_n}{2(1+\alpha_n)}(m_1-m_2)^2
- \frac{m_1^4+m_2^4}{12}
+ O(\|\mathbf{m}\|^5),
\end{align}
%
so that the following bounds are satisfied
\begin{equation}\label{bound-G-crit}
\left\{
\begin{array}{l}
G_{\alpha_n,\beta_c}(\mathbf{m}) 
\ge 2\log 2 -\frac{1}{4}(m_1-m_2)^2 - \frac{m_1^4+m_2^4}{12}
+ O(\|\mathbf{m}\|^5),\\[0.3cm]
G_{\alpha_n,\beta_c}(\mathbf{m}) 
\le 2\log 2 - \frac{m_1^4+m_2^4}{12}
+ O(\|\mathbf{m}\|^5)\,.
\end{array}
\right.
\end{equation}
As for the system outside criticality, we consider, for  $\delta_c \in (1/5, 1/4)$,
the partition of $\Gamma_{n/2}^2$
\begin{equation}\label{def:DNtilde}
\begin{split}
& B_c= \left \{ \bm \in \Gamma_{n/2}^2: \, \|\bm-\bm^0\| \leq \frac{1}{n^{\delta_c}} \right \} \,\, \text{and}
    \,\, \mathcal{\tilde{C}}=\Gamma_{n/2}^2 \setminus B_c \,. 
    \end{split}
\end{equation}
Moreover, we set
\begin{equation}
 \left\{\begin{array}{l}
 R_{\text{low}}^{(n)} =
\mathcal{O}_n B_c,
 \\[0.5cm]
 R_{\text{up}}^{(n)} =
n^{1/4} B_c,
\\
\end{array}
\right.
\qquad
\mbox{ where }
\qquad
\mathcal{O}_n=
\left(
\begin{array}{ll}
\frac{n^{1/4}}{2} & \frac{n^{1/4}}{2}\\
\frac{n^{1/2}}{2} & - \frac{n^{1/2}}{2}
\end{array}
\right)
,
\end{equation}
and define
\begin{equation}
\left\{
\begin{array}{l}
 D_{\text{low}}^{(n)}= 
\sum_{\textbf{x}\in R_{\text{low}}^{(n)}}
\frac{4}{\pi n^{5/4}}
e^{- \frac{1}{2}x_2^2
- \frac{1}{12}x_1^4},
\\[0.5cm]
D_{\text{up}}^{(n)}= 
\sum_{\textbf{x} \in R_{\text{up}}^{(n)}} \frac{4}{\pi n^{3/2}} e^{- \frac{1}{24}(x_1^4+x_2^4)}.
\end{array}
\right.
\end{equation}
Note that by construction 
\begin{equation*}
\left\{
\begin{array}{l}
\lim_{n\to\infty}  D_{\text{low}}^{(n)}= \frac{1}{\pi}\int_{\mathbb{R}^2}
e^{-\frac{1}{2}x_2^2 - \frac{1}{12}x_1^4} dx_1 dx_2 =
\frac{3^{1/4}}{\sqrt{\pi}} \Gamma\!\left(\frac14\right),\\[0.5cm]
\lim_{n\to\infty}  D_{\text{up}}^{(n)}= \frac{1}{4\pi}\int_{\mathbb{R}^2}
e^{-\frac{1}{24}(x_2^4 + x_1^4)} dx_1 dx_2 = \frac{\sqrt6}{8\pi}
\Gamma\!\left(\frac14\right)^2.
\end{array}
\right.
\end{equation*}

In order to prove \eqref{Zcrit-bound}, we start from the trivial observation that
$$
\bar{Z}_{n,\beta_c}(B_c) \le \bar{Z}_{n,\beta_c}=
\bar{Z}_{n,\beta_c}(B_c) +  \bar{Z}_{n,\beta_c}(\mathcal{\tilde{C}})\,.
$$
From the Stirling bound \eqref{StirBound}  and the upper bound on $G_{\alpha_n, \beta_c}$ in \eqref{bound-G-crit}, we get that
\begin{equation}\label{correzione}
    \bar{Z}_{n,\beta_c}(\mathcal{\tilde{C}}) 
     \leq C e^{n \log 2} \sum_{\bm\in \mathcal{\tilde{C}}} e^{-\frac{n}{2}(\frac{m_1^4+m_2^4}{12} + O(\|\bm\|^5))}= o(e^{n \log 2})\,,
    \end{equation}
where in the last step we used that, for all $\bm \in \mathcal{\tilde{C}}$
and since $\delta_c\in (1/5,1/4)$,
\begin{equation*}
n(m_1^4+m_2^4)\ge n^{1-4\delta_c}\ge n^{\tilde{\gamma}}\,, \mbox{ for some }
\tilde{\gamma}>0\,.
\end{equation*}

We now want to show that  $\bar{Z}_{n,\beta_c}(B_c)$ provides 
the correct exponential order $n\log 2$ of the partition function. Starting from the expression of $\bar{Z}_{n,\beta_c}(B_c)$ 
analog to \eqref{eq-Z-tipica}, which can be obtained in $B_c$
using the Stirling formula, and implementing the lower bound in \eqref{bound-G-crit}
and the change of variable ${\bf{x}}=n^{1/4}{\bf{m}}$, we obtain
\begin{equation*}
\begin{split}
    \bar{Z}_{n,\beta_c}(B_c)
    &\le e^{n\log 2} \sum_{\bm \in B_c} \frac{4}{\pi n \sqrt{1-m_1^2} \sqrt{1-m_2^2}}
    e^{-\frac{n}{2}\frac{m_1^4+m_2^4}{12}} (1+o(1))\\
    &=e^{n\log 2}\sum_{{\bf x} \in R_{\text{up}}^{(n)}} \frac{4}{\pi n} e^{-\frac{1}{24}(x_1^4+x_2^4)} (1+o(1)) \\
    &=\sqrt{n}\cdot e^{n\log 2} D_{\text{up}}^{(n)} (1+o(1)),
\end{split}
\end{equation*}
where we used that $n\|\bm\|^5\le n^{1-5\delta_c}=o(1)$. 
Together with \eqref{correzione}, this provides the upper bound in \eqref{Zcrit-bound}.

For the lower bound on  $\bar{Z}_{n,\beta_c}(B_c)$, and hence on $\bar{Z}_{n,\beta_c}$,
 we can analogously implement the upper bound in \eqref{bound-G-crit},  
 and the change of variable $\textbf{x}= \mathcal{O}_n\bm$, or explicitly
$$x_1= \frac{n^{1/4}}{2} (m_1+m_2),\qquad x_2= \frac{n^{1/2}}{2} (m_1-m_2)\,.$$
We then get
\begin{equation*}
\begin{split}
    \bar{Z}_{n,\beta_c}(B_c)
    &\ge e^{n\log 2} \sum_{\bm \in B_c} \frac{4}{\pi n \sqrt{1-m_1^2} \sqrt{1-m_2^2}}e^{-\frac{n}{2}\left(\frac{(m_1-m_2)^2}{2} + \frac{m_1^4+m_2^4}{12}\right)} (1+o(1))\\
    &=e^{n\log 2}\sum_{{\bf x} \in R_{\text{low}}^{(n)}} 
    \frac{4}{\pi n} e^{-\frac{1}{2}x_1^2+ \frac{1}{12}x_2^4} (1+o(1)) \\
    &=\sqrt[4]{n}\cdot e^{n\log 2} D_{\text{low}}^{(n)} (1+o(1)),
\end{split}
\end{equation*}
which concludes the proof of the lower bound in \eqref{Zcrit-bound} and 
of the proposition.
\end{proof}

\end{document}
